% Generated draft: requires mathematical and source review.
\documentclass[11pt]{article}
\usepackage[margin=25mm]{geometry}
\usepackage{fontspec}
\setmainfont{Latin Modern Roman}
\usepackage{amsmath,amssymb,mathtools,unicode-math}
\setmathfont{Latin Modern Math}
\usepackage{enumitem,needspace,xurl,hyperref}
\setlength{\parindent}{0pt}
\setlength{\parskip}{5pt}
\setlength{\emergencystretch}{3em}
\newcommand{\N}{\mathbb{N}}
\newcommand{\Z}{\mathbb{Z}}
\newcommand{\Q}{\mathbb{Q}}
\newcommand{\R}{\mathbb{R}}
\newcommand{\C}{\mathbb{C}}
\newcommand{\F}{\mathbb{F}}
\newcommand{\A}{\mathbb{A}}
\newcommand{\PP}{\mathbb{P}}
\newcommand{\E}{\mathbb{E}}
\newcommand{\eps}{\varepsilon}
\newcommand{\dd}{\,\mathrm{d}}
\begin{document}
\section*{UnsolvedMath ID: 20001646}
\textbf{Source problem code:} catalogue entry
\subsection{Definitions and mathematical statement}
The source's main sentence, ``Does every finite index subgroup contain one of these \(\Gamma\)?'', refers to a family defined before the supplied excerpt. Its precise marked quotient convention is absent. The aid reconstructs a free-group congruence question and separately treats the appended largeness question; these must remain distinct.

For \(F_n\) a free group of finite rank \(n\geq1\) and \(Q\) a finite group, define
\[
 C_Q=\bigcap_{\theta\in\operatorname{Hom}(F_n,Q)}\ker\theta.
\]
The aid reports an effectively computable fully characteristic finite-index subgroup with
\[
 [F_n:C_Q]\leq|Q|^{|Q|^n}.
\]
There are exactly \(|Q|^n\) homomorphisms, determined by generator images, and the diagonal map into their finite product gives the bound. Fully characteristic includes invariance under every endomorphism in the sense \(\alpha(C_Q)\subseteq C_Q\); automorphisms preserve it with equality. Effective computation assumes the finite multiplication table of \(Q\) and a free basis are given.

For an epimorphism \(p:F_n\twoheadrightarrow Q\), define the marked subgroup in the automorphism group by
\[
 \Gamma_p^{\mathrm{Aut}}=\{\alpha\in\operatorname{Aut}(F_n):p\circ\alpha=p\}.
\]
Then the aid's refinement statement is
\[
 \ker(\operatorname{Aut}(F_n)\to\operatorname{Aut}(F_n/C_Q))
 \subseteq\Gamma_p^{\mathrm{Aut}}.
\]
For any characteristic finite-index \(H\leq F_n\), write \(\pi:\operatorname{Aut}(F_n)\to\operatorname{Out}(F_n)\). The reported exact transfer is
\[
 \pi\bigl(\ker(\operatorname{Aut}(F_n)\to\operatorname{Aut}(F_n/H))\bigr)
 =\ker(\operatorname{Out}(F_n)\to\operatorname{Out}(F_n/H)).
\]
An induced inner automorphism of \(F_n/H\) can be removed by lifting its conjugating element to \(F_n\), so the two sides agree. Thus characteristic principal kernels and images of marked quotient subgroups form cofinal congruence systems. This is cofinality within the congruence system; it does not prove that every finite-index subgroup of \(\operatorname{Out}(F_n)\) contains a principal kernel. The aid reports that full CSP question open for \(n\geq3\).

The appended question of Bridson asks, for \(n\geq4\), whether \(\operatorname{Out}(F_n)\) is large: some finite-index subgroup admits an epimorphism onto \(F_2\), equivalently onto a nonabelian free group. The source's parenthetical reference to acting on a tree alone is not equivalent without further hypotheses. The aid reports published property-\((T)\) results ruling out largeness for \(n\geq5\). Since property \((T)\) passes to finite-index subgroups and quotients, while \(F_2\) lacks it, that gives the stated obstruction. For rank four the aid relies on a computer-assisted preprint; its claimed negative conclusion retains that publication and verification qualification and is not promoted here to an independently checked published theorem.

The surrounding source discussion records Bestvina's request for effective constructions, Bridson's suggestions of verbal subgroups or Burnside quotients, and Vogtmann's comment literally referring to characteristic subgroups of \(\operatorname{Out}(F_n)\). The constructed \(C_Q\) are instead subgroups of \(F_n\); no classification of characteristic subgroups of the outer automorphism group is established. The excerpt ends ``On the board, he writes:'' without the ensuing formula, which is not reconstructed.
\subsection{Short English statement}
The aid gives computable characteristic refinements and exact transfer of principal kernels from Aut to Out, but these do not solve full congruence cofinality. The separate largeness question has reported property-\((T)\) obstructions with different published and preprint qualifications.
\subsection{Sources}
\begin{itemize}
\item Local registry: permanent UnsolvedMath ID \(20001646\), exact public code \texttt{catalogue entry}; \texttt{aim-geometric-group-theory-notes.json}, \texttt{notes[154]}.
\item Johanna Mangahas's summary of the AIM workshop ``The geometry of the outer automorphism group of a free group'', October 25--29, 2010, Problem 3.4: \url{https://aimath.org/WWN/outerauto/outerauto.pdf}. The page-4 editor header is extraction spillover. The text names Mladen Bestvina, Martin Bridson and Karen Vogtmann with the contributions preserved above. The literature field is empty.
\item Research attempt 1: \texttt{partial\_result}, status at run \texttt{open}, rigor label \texttt{proved}, novelty confidence low, classification \texttt{PARTIAL-PROGRESS}. It reports the \(C_Q\) bound and Aut--Out transfer, the still-open higher-rank CSP, published property-\((T)\) results for ranks at least five, and a computer-assisted preprint in rank four. No full property-\((T)\) citations or preprint identifier are supplied. The current registry status is partially_solved.
\end{itemize}
\end{document}
